\section{Equivalent crossing problems in diagonal centralizers}
\label{sec:centralizer-gates}

The crossing search can be restricted to full diagonal centralizers
without changing its asymptotic exponent. These are much larger than
the fixed local algebras studied in Section~\ref{sec:hurwitz-amplification}.
The full centralizers retain the unrestricted
asymptotic crossing exponent: a compact-group average incurs only
a polynomial rank overhead. Throughout this section the field is
$\mathbb C$, and $P(q)$ is the unrestricted two-cut cost in
\eqref{eq:pt-crossing-definition}.

\begin{proposition}[Orthogonal averaging with polynomial overhead]
\label{prop:orthogonal-average-gate}
Fix integers $r\ge2$ and $k\ge1$, put $q=r^k$, and identify
$U=(\mathbb C^r)^{\otimes k}$. Let
\[
 \mathcal C_{r,k}=
 \operatorname{End}_{\mathrm O(r,\mathbb R)}(U\otimes U),
\]
where $g$ acts diagonally as $g^{\otimes2k}$. Define
\[
 P_{\mathrm O}(r,k)=
 \min\{\operatorname{rank}M+\operatorname{rank}N:
       I=M+\Gamma N,\ M,N\in\mathcal C_{r,k}\}.
\]
Then
\begin{equation}\label{eq:orthogonal-average-cost}
 P(r^k)\le P_{\mathrm O}(r,k)
 \le\binom{4k+r^2-1}{r^2-1}P(r^k).
\end{equation}
The same bound holds with $\mathrm{SO}(r,\mathbb R)$.
For fixed $r$, both the upper and the lower asymptotic exponents
are unchanged by replacing $P(q)$, for all integer $q$, by
$P_{\mathrm O}(r,k)$ at $q=r^k$.
\end{proposition}

\begin{proof}
Write $R(g)=g^{\otimes2k}$ and let
$L_g(X)=R(g)XR(g)^{\mathsf T}$. Every matrix entry of the
linear operator $L_g$ is a homogeneous polynomial of degree $4k$
in the $r^2$ entries of $g$. Expanding in monomials therefore gives
\[
 \dim\operatorname{span}\{L_g:g\in\mathrm O(r,\mathbb R)\}
 \le D:=\binom{4k+r^2-1}{r^2-1}.
\]
The normalized Haar average
$\mathcal R=\int L_g\,dg$ belongs to this finite-dimensional
span. Indeed, every linear functional annihilating the span also
annihilates the integral. Choosing a basis from the maps $L_g$
expresses $\mathcal R$ as a linear combination of at most $D$
conjugations from the group. Consequently
\[
 \operatorname{rank}\mathcal R(X)\le D\operatorname{rank}X.
\]
This exact finite expression is simultaneous for all $X$; its
coefficients need not be positive.

The average fixes $I$ and has image in $\mathcal C_{r,k}$.
It commutes with $\Gamma$: for $h=g^{\otimes k}$ the elementary
identity $(hYh^{\mathsf T})^{\mathsf T}=hY^{\mathsf T}h^{\mathsf T}$
proves this on simple tensors, hence on all matrices. Averaging
$I=M+\Gamma N$ and applying the rank bound to both summands gives
the upper inequality in \eqref{eq:orthogonal-average-cost}.
The lower inequality follows by restricting the feasible set.
Nothing in the argument changes for the special orthogonal group.

Compression to the first $q$ coordinates commutes with partial
transpose and does not increase rank, so $P(q)$ is nondecreasing.
For $r^k\le q\le r^{k+1}$ one therefore has
\[
 P(r^k)\le P(q)\le P(r^{k+1}).
\]
Divide logarithms by $\log q$, which lies between $k\log r$
and $(k+1)\log r$. Both the limsup and liminf agree along the
geometric subsequence and along all integers. The factor $D$
has logarithm $O_r(\log k)$, so it changes neither exponent.
\end{proof}

The classical orthogonal centralizer theorem identifies
$\mathcal C_{r,k}$ with the image of the Brauer algebra on $2k$
strands at parameter $r$. We use this as a cited external
identification; see \cite[Section~3.2 and Corollary~3.8]{LZ2012}
for the action by place permutations and pair contractions and its
surjectivity, attributed there to Brauer. The compact real and
complex orthogonal centralizers coincide by complexification.
At small $r$ the word ``image'' is essential: no faithfulness or
semisimplicity of the abstract Brauer algebra is assumed.
This centralizer contains contractions and permutations coupling
different copies. It is not the commutative product algebra
$\mathcal B_r^{\otimes k}$.

\begin{corollary}[A diagonal-centralizer square gate]
\label{cor:orthogonal-price-gate}
Under the hypotheses of Theorem~\ref{thm:square-pt-exponent},
the Euclidean square is not a price if and only if there exist
$\varepsilon>0$ and $k_0$ such that
\[
 P_{\mathrm O}(4,k)\le4^{(\sqrt2-\varepsilon)k}
 \qquad(k\ge k_0).
\]
\end{corollary}

\begin{proof}
An eventual power improvement for the restricted cost is one for
$P(4^k)$, hence for all dimensions by the monotonicity argument.
Conversely an unrestricted improvement absorbs the polynomial
factor in \eqref{eq:orthogonal-average-cost} after reducing the
positive exponent slack. Apply Theorem~\ref{thm:square-pt-exponent}.
\end{proof}

\subsection{A spinor gate in the Temperley--Lieb centralizer}

The two-dimensional alternating form gives a smaller polynomial
overhead and a description by noncrossing diagrams. The transpose
convention matters: it changes the second summand by a rank-preserving
conjugation, and does not identify the corresponding self equations.

\begin{proposition}[Temperley--Lieb reduction of the unrestricted cost]
\label{prop:temperley-lieb-gate}
Let $W=\mathbb C^2$, $U=W^{\otimes k}$, $q=2^k$, and put
\[
 \epsilon=\begin{pmatrix}0&1\\-1&0\end{pmatrix},\qquad
 J=\epsilon^{\otimes k},\qquad
 t_J(X)=J^{-1}X^{\mathsf T}J,\qquad
 \Gamma_J=\operatorname{id}\otimes t_J.
\]
Let $\mathcal T_k=\operatorname{End}_{\mathrm{SL}_2(\mathbb C)}
(W^{\otimes2k})$ for the diagonal action, and define
\[
 P_{\mathrm{TL}}(k)=
 \min\{\operatorname{rank}M+\operatorname{rank}N:
 I=M+\Gamma_JN,\ M,N\in\mathcal T_k\}.
\]
For every $k\ge1$,
\begin{equation}\label{eq:temperley-lieb-overhead}
 P(2^k)\le P_{\mathrm{TL}}(k)
       \le\binom{4k+3}{3}P(2^k).
\end{equation}
The algebra $\mathcal T_k$ is the image of
$\mathrm{TL}_{2k}(2)$, generated by $e_i=I-F_i$, where $F_i$
exchanges adjacent spinor factors. It is also the image of the
symmetric-group algebra on these factors, and $\Gamma_J$ preserves
it. Both asymptotic exponents of $P$ are unchanged by this reduction.
\end{proposition}

\begin{proof}
The matrix $J$ is real orthogonal, so $J^{-\mathsf T}=J$.
If $\mathcal A_J$ denotes conjugation by $I\otimes J$, then
\begin{equation}\label{eq:spinor-transpose-conversion}
 \Gamma_J\mathcal A_J(X)=\Gamma(X).
\end{equation}
To verify it, apply the two sides to a simple tensor and use
\[
 J^{-1}(JYJ^{-1})^{\mathsf T}J
  =J^{-1}J^{-\mathsf T}Y^{\mathsf T}J^{\mathsf T}J
  =Y^{\mathsf T}.
\]
Thus conjugating only the second summand gives a bijection between
unrestricted decompositions for $\Gamma$ and $\Gamma_J$, preserving
both ranks. This holds also for odd $k$, when $J$ is alternating.

Every $g\in\mathrm{SL}_2(\mathbb C)$ preserves $\epsilon$.
For $h=g^{\otimes k}$ this gives $h^{\mathsf T}Jh=J$, hence
\[
 t_J(hYh^{-1})=h\,t_J(Y)\,h^{-1}.
\]
It follows that $\Gamma_J$ commutes with conjugation by
$h\otimes h=g^{\otimes2k}$. Average these conjugations over
$\mathrm{SU}(2)$. The average fixes $I$ and lands in $\mathcal T_k$:
the compact and complex centralizers agree because the complex
span of $\mathfrak{su}(2)$ is $\mathfrak{sl}_2(\mathbb C)$.

The entries of $g^{\otimes2k}$ are homogeneous of degree $2k$
in the four entries of $g$. Since $\det g=1$, the inverse of $g$
is its adjugate and is linear in those same entries. Therefore
the conjugation superoperator has homogeneous degree $4k$.
Its span has dimension at most $\binom{4k+3}{3}$. The exact
finite-span averaging argument in
Proposition~\ref{prop:orthogonal-average-gate} now bounds the
rank of each averaged summand by this factor times its original
rank. Together with \eqref{eq:spinor-transpose-conversion} this
proves \eqref{eq:temperley-lieb-overhead}.

The centralizer identification is classical dimension-two
Schur--Weyl duality. A precise primary reference is
\cite[Section~6.1, Theorem~6.4, pp.~28--29]{CruzErdmann2026}, whose convention
uses $U_i=F_i-I$ and parameter $-2$; taking $e_i=-U_i$ gives
parameter $2$. Equivalently, the adjacent permutations generate
the commutant and the antisymmetrizer on three spinors vanishes.
The resulting Temperley--Lieb action has its usual noncrossing
pairing basis. The covariance proved above shows that $\Gamma_J$
preserves the commutant. Finally, the polynomial overhead and
monotonicity of $P(q)$ give both exponent equalities exactly as in
Proposition~\ref{prop:orthogonal-average-gate}.
\end{proof}

\begin{corollary}[A noncrossing-diagram square gate]
\label{cor:temperley-lieb-price-gate}
Under the hypotheses of Theorem~\ref{thm:square-pt-exponent},
the Euclidean square is not a price if and only if there exist
$\varepsilon>0$ and $k_0$ such that
\[
 P_{\mathrm{TL}}(k)\le2^{(\sqrt2-\varepsilon)k}
 \qquad(k\ge k_0).
\]
\end{corollary}

\begin{proof}
Apply \eqref{eq:temperley-lieb-overhead}, absorb its polynomial
factor by reducing the positive exponent slack, and invoke
Theorem~\ref{thm:square-pt-exponent}.
\end{proof}

The permitted operators here are arbitrary linear combinations
of the noncrossing pairing operators on all $2k$ spinor strands.
They couple different copies and form a noncommutative algebra.
Consequently the bounds below for tensor products of local
flip--cup or Hodge algebras do not decide this gate. No rank bound
for arbitrary planar combinations, or multiplicative rule for
the minimum two-cut cost, is asserted.

\subsection{A fixed octonion output span cannot be split more cheaply}

The classical eight-square composition formula gives a two-cut
cost of sixteen at input dimension eight. One possible finite
improvement would divide its output space between the two cuts.
The following calculation excludes that particular operation.
It imposes no restriction on other output spaces.

\begin{proposition}[Rigidity of the fixed octonion output span]
\label{prop:octonion-fixed-span}
Let $\kk$ have characteristic different from two. Form the
quaternion algebra with basis $1,i,j,k$, relations
$i^2=j^2=k^2=ijk=-1$, and its usual conjugation. On pairs of
quaternions use the Cayley--Dickson product
\[
 (x,y)(u,v)=(xu-\bar v y,\;vx+y\bar u).
\]
Let $e_0,\ldots,e_7$ be its ordered basis
$(1,i,j,k)$ in the first component followed by the same basis
in the second component. Let $V$ be the $64$ by $8$ matrix
whose row $(a,b)$ is the coefficient vector of $e_ae_b$.
Write $A_R=VRV^{\mathsf T}$ for $R\in\operatorname{Mat}_8(\kk)$.
Both linear maps
\[
 R\longmapsto (\mathrm{Id}+\Gamma)A_R,
 \qquad R\longmapsto(\mathrm{Id}-\Gamma)A_R
\]
are injective. Consequently the only solution of
\[
 I_{64}=A_R+\Gamma A_S
\]
in this fixed common output span is $R=S=I_8/2$.
Its two summands each have rank eight, so its total cost is sixteen.
\end{proposition}
\begin{proof}
The quaternion norm is $N(x)=x\bar x$, the sum of the four
coordinate squares. Expanding the displayed Cayley--Dickson
product gives
\[
 N(xu-\bar v y)+N(vx+y\bar u)
  =(N(x)+N(y))(N(u)+N(v)).
\]
Indeed the unmixed terms give the right-hand side; the two cross
terms cancel because the scalar parts of $xu\bar yv$ and
$vxu\bar y$ agree by cyclicity in the associative quaternion
algebra. Polarizing in the two input blocks therefore gives
\begin{equation}\label{eq:octonion-gram}
 VV^{\mathsf T}+\Gamma(VV^{\mathsf T})=2I_{64}.
\end{equation}

Suppose $A_R+\sigma\Gamma A_R=0$, where $\sigma\in\{1,-1\}$,
and put $r(e_a)=R_{0a}$. The matrix entry with row $(0,0)$
and column $(a,b)$ says
\[
 r(e_ae_b)+\sigma R_{ba}=0.
\]
If $\sigma=1$, taking $b=0$ shows $2r(e_a)=0$ for every $a$,
and then every entry of $R$ is zero. If $\sigma=-1$, the equation
instead gives $R_{ba}=r(e_ae_b)$. The entry with row $(0,j)$
and column $(a,b)$ now gives
\[
 r\bigl((e_ae_b)e_j-(e_ae_j)e_b\bigr)=0.
\]
The following products, computed directly from the defining
quaternion relations, span the eight-dimensional output space:
\[
\begin{array}{c|c@{\qquad}c|c}
(j,a,b)&\text{product difference}&(j,a,b)&\text{product difference}\\ \hline
(1,2,3)&-2e_0&(1,0,5)&2e_4\\
(1,2,2)&-2e_1&(1,0,4)&-2e_5\\
(1,0,3)&2e_2&(1,0,7)&-2e_6\\
(1,0,2)&-2e_3&(1,0,6)&2e_7
\end{array}
\]
Thus $r=0$ and $R=0$, proving both injectivity assertions.

Apply $\mathrm{Id}-\Gamma$ to $I=A_R+\Gamma A_S$.
The second injectivity assertion gives $R=S$.
Equation~\eqref{eq:octonion-gram} and the first injectivity assertion
then give $R=S=I_8/2$. Finally $V^{\mathsf T}V=8I_8$:
each output basis element occurs with a sign in exactly eight
multiplication-table entries, and each entry has one output.
Since eight is invertible, $VRV^{\mathsf T}$ has rank eight for
$R=I_8/2$. This proves the stated cost and its precise scope.
\end{proof}

The proof concerns the $64$-dimensional space of matrices
$VRV^{\mathsf T}$. It does not show that sixteen is the unrestricted
minimum at dimension eight, and it gives no tensor-power rule for
that minimum.
\subsection{A limitation of linear averaging in characteristic two}

The complex averaging proof uses a projection onto invariant
operators. The corresponding linear projection is unavailable in
characteristic two, even before a rank bound is imposed.

\begin{proposition}[Equivariant linear maps kill the identity]
\label{prop:char2-reynolds-obstruction}
Let $\kk$ have characteristic two, let $m\geq1$, and let
$V=(\kk^2)^{\otimes m}$ carry the diagonal action of
$G=\mathrm{SL}_2(\kk)$. Put $Z=\operatorname{End}_G(V)$.
Every $G$-equivariant linear map
\[
 \pi:\operatorname{End}_{\kk}(V)\longrightarrow Z
\]
for the conjugation actions satisfies $\pi(I_V)=0$.
In particular there is no equivariant linear retraction onto $Z$.
The assertion also applies to the algebraic group $\mathrm{SL}_2$.
\end{proposition}
\begin{proof}
Let $P=E_{00}\otimes I$ project onto the first basis vector in
the first spinor factor. In characteristic two,
$s=\left(\begin{smallmatrix}0&1\\1&0\end{smallmatrix}\right)$
belongs to $G$. For $S=s^{\otimes m}$ one has
$SPS^{-1}=I_V-P$. Equivariance and the fact that every value of
$\pi$ commutes with $S$ give
\[
 \pi(I_V)=\pi(P)+\pi(SPS^{-1})
          =\pi(P)+S\pi(P)S^{-1}=2\pi(P)=0.
\]
Since $I_V\in Z$ is nonzero, a retraction cannot have this property.
\end{proof}

This rules out a Reynolds-type equivariant linear averaging map.
It does not rule out a polynomial-rank reduction to the modular
centralizer by a different procedure, or a crossing construction
inside that centralizer.

\subsection{A product-Bell diagonal construction problem}

A second explicit ansatz uses a complete diagonal algebra. It need
not be invariant under the diagonal spinor action, so it is a
separate construction family from the full centralizer above.

\begin{proposition}[Product-Bell diagonal reduction]
\label{prop:bell-diagonal-pt}
Put $q=2^n$, $U=(\mathbb C^2)^{\otimes n}$, and
$J=\epsilon^{\otimes n}$ as above. Pair the corresponding factors
of the two copies of $U$. On each pair let $P_0,P_1,P_2,P_3$ be
the rank-one projectors onto the four standard Bell vectors
\[
 (|00\rangle+|11\rangle)/\sqrt2,\quad
 (|00\rangle-|11\rangle)/\sqrt2,\quad
 (|01\rangle+|10\rangle)/\sqrt2,\quad
 (|01\rangle-|10\rangle)/\sqrt2.
\]
For $a\in\mathbb C^{\{0,1,2,3\}^n}$ set
$M(a)=\sum_xa_xP_{x_1}\otimes\cdots\otimes P_{x_n}$.
With $\mathbf1$ denoting the all-one vector and
\[
 G=\frac12(\mathbf1\mathbf1^{\mathsf T}-2I_4),
\]
one has
\[
 \Gamma_J M(b)=M(G^{\otimes n}b),\qquad
 \operatorname{rank}M(a)=|\operatorname{supp}a|.
\]
Consequently every solution
\begin{equation}\label{eq:bell-diagonal-affine}
 a+G^{\otimes n}b=\mathbf1
\end{equation}
gives an ordinary partial-transpose decomposition of total cost
$|\operatorname{supp}a|+|\operatorname{supp}b|$.
Conversely every decomposition with both twisted summands in this
diagonal algebra is represented by this equation.
\end{proposition}
\begin{proof}
For a two-by-two matrix,
$\epsilon^{-1}X^{\mathsf T}\epsilon=\operatorname{tr}(X)I-X$.
Each Bell projector has partial trace $I_2/2$, and therefore
\[
 \Gamma_\epsilon(P_j)=I_4/2-P_j
       =\sum_iG_{ij}P_i.
\]
Tensoring proves the transform. The mutually orthogonal projectors
have rank one, so the rank is precisely the number of nonzero
eigenvalues. Equation~\eqref{eq:spinor-transpose-conversion}
converts the second twisted summand to an ordinary one by
conjugation, preserving its rank. No restriction on the parity of
$n$ is needed.
\end{proof}

There is also a signed Fourier description. Identify each coordinate
with $\mathbb F_2^2$, put
$Q_0(z)=z_1+z_2+z_1z_2$, and define
$D_{xx}=(-1)^{\sum_iQ_0(x_i)}$.
Let $\mathcal F$ be the normalized Walsh transform for the
coordinatewise symplectic pairing:
\[
 \mathcal F_{xy}=2^{-n}
 (-1)^{\sum_i(x_{i,1}y_{i,2}+x_{i,2}y_{i,1})}.
\]
The quadratic identity
$Q_0(x+y)=Q_0(x)+Q_0(y)+x_1y_2+x_2y_1$
gives
\[
 G^{\otimes n}=(-1)^nD\mathcal F D.
\]
Thus \eqref{eq:bell-diagonal-affine} asks for
\[
 Da+(-1)^n\mathcal F(Db)=D\mathbf1
\]
with small supports on both sides. This is an exact reformulation
of this diagonal ansatz; no claim about the support minimum is
implicit in it.

\begin{example}[A finite dimension-eight witness]
\label{ex:bell-diagonal-eight}
Use base-four integer labels for three coordinates and put
\[
 \begin{split}
 A&=\{0,10,12,20,25,27,38,54\},\\
 B&=\{5,7,10,18,20,30,40,56\}.
 \end{split}
\]
Then $a=4\mathbf1_A$, $b=4\mathbf1_B$ solve
\eqref{eq:bell-diagonal-affine}. Indeed, if
$H=(\mathbf1\mathbf1^{\mathsf T}-2I_4)^{\otimes3}$, direct
integer summation gives
\[
 H\mathbf1_B=2\mathbf1-8\mathbf1_A.
\]
There are only $64$ entries to check: the signed sum is $-6$
on $A$ and $2$ elsewhere. The two ordinary summands consequently
have ranks eight and eight. This recovers a cost-sixteen upper
bound at $q=8$ within an explicitly diagonal family.
It neither proves optimality at $q=8$ nor gives a valid
tensor-power amplification rule.
\end{example}

